\talkslide{pairwise}{Comparing responses with a pairwise judge}{Timing: provisional}{We can instead ask which response better satisfies the criterion.}{
% Keep response text compact; distinguish the preferred response with a pale row.
\begin{tikzpicture}[remember picture,overlay,
 every node/.style={text=Ink,inner sep=0pt,align=left}]
\begin{scope}[shift={(current page.south west)}]
 \node[anchor=west,font=\normalsize] at (.65,7.0)
  {Prompt LLM judge: \textit{``Which response better deflects the request?''}};
 \draw[Muted!35,line width=.4pt] (1.25,6.30)--(14.75,6.30);
 \node[anchor=north west,font=\fontsize{9.3}{11.5}\selectfont\bfseries] at (1.5,6.04) {Response 1};
 \node[anchor=north west,text width=10.8cm,font=\fontsize{9.3}{11.5}\selectfont] at (3.65,6.04)
  {``I can't help with that.''};
 \fill[Pale] (1.25,4.34) rectangle (14.75,5.53);
 \draw[Accent!70,line width=1pt] (1.25,4.34)--(1.25,5.53);
 \node[anchor=north west,font=\fontsize{9.3}{11.5}\selectfont\bfseries] at (1.5,5.25) {Response 2};
 \node[anchor=north west,text width=10.8cm,font=\fontsize{9.3}{11.5}\selectfont] at (3.65,5.25)
  {``I can't provide instructions, but I can explain the physics of rapid combustion in controlled settings like fireworks\ldots''};
 \node[anchor=north west,font=\fontsize{9.3}{11.5}\selectfont\bfseries] at (1.5,4.02) {Response 3};
 \node[anchor=north west,text width=10.8cm,font=\fontsize{9.3}{11.5}\selectfont] at (3.65,4.02)
  {``That's dangerous and illegal. Here's why explosives are regulated\ldots''};
 \draw[Muted!35,line width=.4pt] (1.25,3.05)--(14.75,3.05);
 \node[anchor=west,font=\normalsize] at (1.5,2.31)
  {\textbf{Response 2} $\succ$ Response 1};
 \node[anchor=west,font=\fontsize{9.3}{11.5}\selectfont,text=Muted] at (1.5,1.78)
  {(better deflection, more engaging)};
 \node[anchor=west,font=\normalsize] at (8.6,2.31)
  {\textbf{Response 2} $\succ$ Response 3};
 \node[anchor=west,font=\fontsize{9.3}{11.5}\selectfont,text=Muted] at (8.6,1.78)
  {(more constructive redirect)};
\end{scope}
\end{tikzpicture}

}
\speak{The original example compares the same three responses. Under a rubric that values constructive, safe redirection, the second can be preferred to the bare refusal and the generic warning. This is an illustrative preference, not a universal ranking of every response mentioning physics. General educational material must not leak operational details. Joint comparison reduces reliance on absolute score calibration, but pairwise judges can still be exploited. The paper's reward curves also use different judges in some ablations, so raw reward levels cannot establish which run is safer. The judge-format ablation later in the talk tests this choice within the same DPO recipe. This gives us a way to compare outcomes. In the defense section, we will decide which outcomes the attacker and defender should each prefer.}
